\documentclass[11pt]{article}
\usepackage[margin=1in]{geometry}
\usepackage[T1]{fontenc}
\usepackage[utf8]{inputenc}
\usepackage{lmodern}
\usepackage{microtype}
\usepackage{booktabs}
\usepackage{url}
\usepackage[hidelinks]{hyperref}
\usepackage{xcolor}

\title{RIGOR: Research Integrity Guardrails for Open Research\\
\large An integrity-first agent for the full research workflow ---\\
grounded literature, verified experiments, honest statistics, auditable manuscripts}
\author{Alif Mahmud\\
\small Department of Electrical and Electronic Engineering\\
\small Bangladesh University of Engineering and Technology (BUET), Dhaka, Bangladesh}
\date{September 2026}

\begin{document}
\maketitle

\begin{abstract}
Large language model (LLM) agents are moving into every stage of the research
workflow---literature survey, experiment execution, statistical analysis, and
manuscript drafting---and they import three well-documented failure modes:
hallucinated citations, numbers in prose that silently drift from the data
behind them, and superiority claims that random seed variation does not
support. We present RIGOR, an open-source toolkit of sixteen agent
\emph{skills} (instruction files paired with small, dependency-light Python
programs) that makes agent-assisted research \emph{grounded by construction}
rather than by exhortation. RIGOR enforces three architectural guarantees:
(1)~a paper may be cited only if it exists in a provenance-tagged store
populated exclusively by live bibliographic API responses; (2)~audit tools
that reconcile bibliographies and manuscript numbers against ground truth
\emph{report and propose but never edit}; and (3)~every ``X beats Y'' claim is
gated on paired-by-seed significance tests that report exact $p$-values and
$n$, with non-significance reported as a result. Around these gates, RIGOR
automates the experiment loop itself: the agent writes and edits the
experiment code, proves every pipeline on a local smoke subset before any
GPU is spent, dispatches multi-seed sweeps to free cloud GPUs (Kaggle
unattended and recommended; Colab semi-attended and field-verified, where one
human tap doubles as the GPU
approval)---optionally scheduling across a research team's pooled per-account
quotas---and passes every completed run through an integrity checklist
before its numbers reach the researcher. The manuscript is treated the same
way: \LaTeX{} compiled locally, diagrams written as TikZ code the agent
renders and visually inspects itself, and every result table generated from
raw run artifacts---numbers flow one way, from code to prose. Investigations
that outlive a session are carried by an append-only lab notebook whose
findings cite evidence artifacts and whose agent-written narrative is
mechanically checked against the log. Around the submission itself:
equivalence claims are gated on TOST rather than inferred from failed
difference tests; retracted and miscited references are flagged; datasets are
fingerprinted against silent drift; and a one-command pre-submission gate,
SHA-256 artifact freezing, and diff-verified reviewer responses close the
loop. The same gates also run reviewer-side: two skills screen submitted PDFs
for references that resolve to no existing work and for numbers a paper
contradicts itself on, emitting quoted candidates that must pass human
adjudication before reaching a program chair. On a 314-paper conference
screening batch, adjudication confirmed 26 untraceable references and 15
internal inconsistencies---and cleared most raw flags as legitimate, which is
the case for making adjudication mandatory. We describe the
design, a
two-layer portability model that separates generic mechanics from per-project
configuration, and a working case study on a real multi-seed forecasting
project in which each tool produced verifiable findings, including a wrong
bibliography title, 47 stale-drift candidates in one manuscript, and a
single-seed ``best model'' headline that dissolved under paired testing.
RIGOR is available under the MIT license.
\end{abstract}

\section{Introduction}

LLM assistance in research is no longer hypothetical: models survey fields,
draft related-work sections, orchestrate experiments, and edit manuscripts.
The efficiency gains are real, but so are the integrity risks, and they are
not hypothetical either. Controlled evaluations find that a large fraction of
LLM-generated bibliographic citations are fabricated or substantively
erroneous~\cite{walters2023fabrication}. Two further failure modes emerge as
agents move beyond writing into the experimental loop:

\begin{enumerate}
\item \textbf{Text--data drift.} Modern empirical papers regenerate their
result tables from experiment artifacts, but the surrounding prose---the
abstract, discussion sentences, figure captions---is hand-written. After a
re-run, the tables update and the prose silently keeps the old numbers. In
our case study (Section~\ref{sec:case}), a careful \emph{manual} audit of one
manuscript found nine such mismatches; the mechanized audit later surfaced
dozens of candidates.
\item \textbf{Seed-noise claims.} With $n{\approx}10$ random seeds, a nominal
gap between two models is often smaller than the seed-to-seed spread of either
one. A single-seed ranking presented as a finding is not robust, and agents
that summarize results tables reproduce such rankings uncritically.
\end{enumerate}

Prompting an agent to ``be accurate'' does not fix any of these: the failure
happens precisely when the model is confident. Our position is architectural:
\emph{integrity failures should be surfaced by mechanism, not caught by
vigilance}. RIGOR contributes a small, auditable toolkit organized around
three guarantees:

\begin{itemize}
\item \textbf{Grounded by construction} (Section~\ref{sec:lit}): every
citation must originate from a live Semantic Scholar API
response~\cite{kinney2023semanticscholar} recorded in a provenance-tagged
store; remembered papers must be verified against the API before they may be
mentioned; claims about a paper's content must trace to its stored abstract,
its full text (with page tags), or the citing sentences retrieved for it.
\item \textbf{Report, propose, never edit} (Section~\ref{sec:audit}): the
bibliography and manuscript auditors emit classified findings with suggested
fixes; a human applies each one. There is no code path by which an audit
modifies the user's files.
\item \textbf{Exact statistics} (Section~\ref{sec:stat}): superiority claims
are tested paired-by-seed (Wilcoxon signed-rank~\cite{wilcoxon1945individual}
with a paired $t$ reference~\cite{virtanen2020scipy}, optional Holm
correction~\cite{holm1979simple}); output always includes $n$ and exact
$p$-values, and a null result is reported, never dropped.
\end{itemize}

These guarantees would matter little if the agent merely \emph{advised} while
a human ran everything by hand. RIGOR therefore also automates the loop the
guarantees protect (Section~\ref{sec:auto}): the agent prepares and edits the
experiment code itself, proves it end-to-end on a local smoke subset,
dispatches unattended multi-seed sweeps to free cloud GPUs, and verifies
everything that comes back---so the integrity gates sit inside a workflow
that is genuinely hands-off between the researcher's kickoff instruction and
the returned briefing.

The guarantees are also symmetric. A program committee meets the same failure
modes from the other direction---references that point at nothing, headline
numbers the paper's own tables contradict---with dozens of submissions and no
time to check either by hand. Two skills therefore run RIGOR's reference and
consistency audits on submitted PDFs, with mandatory human adjudication before
anything reaches a chair (Section~\ref{sec:review}), and were exercised on a
real conference batch before being claimed here (Section~\ref{sec:field}).

\section{Related work}

Commercial AI research assistants (literature-discovery products and chat
interfaces with retrieval) address parts of the citation-grounding problem,
but are generally closed-source, non-scriptable, and difficult to audit; their
guarantees, where stated, are not inspectable. On the open side,
bibliographic APIs---Semantic Scholar~\cite{kinney2023semanticscholar},
OpenAlex~\cite{priem2022openalex}, and Crossref~\cite{hendricks2020crossref}%
---provide the verifiable substrate RIGOR builds on, but are typically consumed
by ad-hoc scripts without an integrity policy connecting them to what the
agent ultimately writes. Reproducibility tooling (workflow managers,
experiment trackers) addresses artifact provenance but not the last mile where
numbers enter prose. Open manuscript-preparation systems come closest in
spirit: Manubot assembles papers from version-controlled sources with
citation-by-identifier~\cite{himmelstein2019manubot}, and showyourwork ties an
article's figures to the code that regenerates
them~\cite{showyourwork}. Both enforce reproducible \emph{manuscript assembly};
RIGOR is complementary, governing the agent-driven loop that produces the
numbers and citations \emph{before} they reach the manuscript. RIGOR's
contribution is not a new algorithm but the missing policy layer: a set of
enforceable rules, each backed by a mechanical tool, spanning literature,
bibliography, manuscript, and statistics.

\section{Design}

\subsection{Skills: judgment and mechanics, separately}

Each RIGOR capability is a \emph{skill}: a folder containing a
\texttt{SKILL.md} (the agent-facing contract: when to act, integrity rules,
how to adjudicate tool output) and---in all but the remote-execution skill,
which drives a shared runner template---one Python script (the mechanics:
API calls, parsing, classification). This split is deliberate. The scripts are
deterministic and testable offline (all HTTP is stubbed in CI); the judgment
calls---is this near-miss number a stale claim or a coincidentally similar
correlation coefficient?---are documented duties of the agent, executed under
human review. The scripts are standard-library-first Python~3.10+ (scipy and
pypdf~\cite{pypdf} only where noted) and equally usable by hand, with no agent
at all.

\subsection{Two-layer portability}

Skills contain no project paths, project names, or reference numbers
(\emph{Layer~1}). Everything project-specific lives in one profile file
(\emph{Layer~2}): manuscript path, results glob, which interpreter has scipy,
where verified reference numbers live, and safety flags. The agent reads the
profile and wires values into command-line arguments; the scripts stay
generic. Installing RIGOR into a project is a folder copy; onboarding scans
the repository, asks only what cannot be inferred, and writes the profile.
Credentials (API keys, cloud tokens) are machine-level environment variables
and never travel with the folder.

\subsection{Grounded literature work}
\label{sec:lit}

The \textbf{lit-review} skill maintains, per topic, a JSON store of papers
keyed by Semantic Scholar identifier, where every entry records
\emph{provenance}: which query, snowball, recommendation, or lookup produced
it. Sub-commands cover ranked search; exhaustive boolean \emph{bulk} search
for long-tail coverage claims; citation snowballing in both directions;
recommendations; \emph{contexts} (the sentences in which the field cites a
paper, with intent tags--- upgrading novelty screening from
regex-over-abstracts to reading actual usage); \emph{enrich} (abstracts
missing from Semantic Scholar are filled from
OpenAlex~\cite{priem2022openalex} into a separately-attributed field, never
overwriting the primary source); \emph{refresh} (batch updates of citation
counts, stamping every entry with a fetch date so ``cites $N$'' claims carry
an as-of date); validated open-access PDF retrieval (magic-byte checks with an
arXiv fallback, because publisher ``OA'' links are sometimes HTML landing
pages); and page-tagged full-text extraction, so a claim grounded in a paper's
body cites a page.

The \emph{report} command renders the collection as a document a researcher
actually reads: every paper appears as a rich entry---title, Semantic
Scholar / DOI / open-access-PDF links, citation counts, provenance, and the
best available abstract with its source attributed---and, given the
researcher's own project description as a \emph{focus}, the whole report is
re-ranked and tiered (\textsc{core} / \textsc{related} /
\textsc{peripheral}) by relevance to that focus. Relevance is a transparent
IDF-weighted term-overlap score (title weighted over summary over abstract)
whose matched terms are printed beside every paper, so the ranking is
inspectable and correctable by eye rather than a black box. The integrity
rules are short and absolute: no paper may be cited unless it is in the
store; remembered papers are verified first; content claims trace to stored
fields; citation counts are reported exactly as returned.
\textbf{topic-watch} re-runs the store's own recorded queries and diffs for
new papers, keeping a survey current before a revision.

The same grounding extends from papers to people. \textbf{pi-scout} helps an
early-career researcher find prospective supervisors by starting from their
\emph{own} papers rather than a generic topic search: it harvests recent work
that cites or relates to those papers, fetches the author pool's metrics in
one batch, and ranks likely principal investigators with a transparent
additive score whose components---citing the researcher's work, last-author
roles, breadth, $h$-index band, recency---are printed per candidate. The
script stops where the API's facts stop. Current position, institutional rank,
and funding are verified by the agent on official pages, each with a source
URL and access date; an $h$-index is never estimated, a job title is never
inferred from a stale affiliation string, and the tool contacts no one.

\subsection{Auditing the bibliography and the manuscript}
\label{sec:audit}

\textbf{bib-audit} checks each BibTeX entry against Semantic Scholar and
Crossref: DOI-first, title-match fallback with similarity scoring. Verdicts
are \textsc{verified}, \textsc{mismatch} (field-level diffs and suggested
corrections), \textsc{not-found}, \textsc{non-paper-ok} (datasets, standards,
software: checked by URL liveness instead), and \textsc{unverifiable}
(submitted work). Two calibration details matter in practice. First,
bibliographic databases often record a conference paper under its preprint
year; a $\le$2-year offset is therefore reported as an informational note, not
an error---na\"ively flagging it would have produced false alarms on a third
of our case-study bibliography. Second, when a title-only match disagrees on
year by more than two years, the finding is flagged but \emph{no fix is
auto-suggested}: the match is likely a different edition or a review of the
same work, and a plausible-looking auto-fix would be worse than none. Each
confidently-identified work is additionally checked against Retraction Watch
data (via OpenAlex's \texttt{is\_retracted}); a \textsc{retracted} verdict
ranks above every other finding, with the instruction to remove the citation
or cite the work explicitly \emph{as} retracted.

\textbf{cite-check} closes the gap existence checks cannot: a real paper cited
for a claim it never makes. The script extracts every citation-bearing
sentence, resolves each key against the bibliography, and pairs it with the
cited work's abstract---drawn only from lit-review stores, so provenance
holds and the tool itself needs no network. The agent then judges each pairing
\textsc{supported} (only with a verbatim quote from the abstract---no quote,
no verdict), \textsc{not-supported} (a candidate for the human, never an
auto-edit), or \textsc{cant-verify} (the claim lives outside what an abstract
covers; the fulltext path is suggested). Judging from model memory of a paper
is prohibited: if the evidence is not in the worksheet, the verdict is
\textsc{cant-verify}, not a recollection.

\textbf{claims-audit} extracts every numeric literal from a LaTeX or Markdown
manuscript's prose and captions (skipping generated table inputs, references,
labels, layout dimensions---and, for Markdown, frontmatter, code blocks,
citation keys, and URLs) and classifies each against a ground-truth pool built from
the machine-generated tables plus the aggregated experiment results:
\textsc{matched} (within rounding), \textsc{near-miss} (close but drifted:
the prime stale-claim suspects), and \textsc{orphan} (traceable to no data
source---mostly legitimate derived or structural values, but containing the
dangerous class of results-like numbers with no provenance). It also flags
unfilled placeholder macros and any results-derived figure whose file is older
than the newest results artifact. The tool deliberately trades false positives
for zero false negatives on the dangerous class; the human pass over the
report is the point, not a failure of automation.

\subsection{Statistical honesty}
\label{sec:stat}

\textbf{stat-check} ingests per-run results files (one JSON per experiment
run), groups them into studies, deduplicates re-runs (newest per seed wins),
excludes smoke tests, and aligns metric values \emph{by seed} so tests are
properly paired. For each requested pair it reports $n$, means, the median
paired difference, the exact Wilcoxon signed-rank $p$, and a paired-$t$
reference value; a default battery tests best-vs-runner-up and best-vs-each-%
baseline. With $n{=}10$ seeds the smallest achievable two-sided Wilcoxon $p$
is $2/2^{10} \approx 0.002$; the tool documents this floor so users read
$p=0.002$ as ``a clean sweep,'' not as arbitrary precision. Holm correction is
available and its use (or absence) is disclosed in the output.

Each pair also carries a 95\% confidence interval on the mean paired
difference and Cohen's $d_z$, so ``significant'' can be qualified by ``how
big.'' Most importantly, the tool enforces the asymmetry that informal
practice erases: a failed difference test is \emph{not} evidence of
equivalence. Claiming two configurations are ``equivalent,'' ``tied,'' or
``statistically indistinguishable'' requires a TOST equivalence test
(\texttt{--tost}, two one-sided paired $t$ against a pre-declared margin);
a pair can be neither significantly different nor equivalent, and the honest
verdict for that state is \textsc{inconclusive}.

\subsection{Data provenance: the bug you haven't noticed yet}
\label{sec:data}

The most expensive research bugs are not crashes but datasets that silently
stop meaning what the researcher thinks they mean---a categorical feature
column that a merge-key type mismatch turned all-constant, a stale artifact
averaged into a multi-seed mean. Nothing fails; the numbers are just quietly
wrong. \textbf{data-audit} makes this class mechanical: \texttt{fingerprint}
records, per data file, the shape, per-column null counts, distinct counts,
ranges, and content hashes, and immediately screams on degeneracies
(all-null, majority-null, or constant columns; byte-identical duplicate
columns; NaN-bearing arrays)---a first fingerprint of already-broken data
must shout, not just record. \texttt{verify} recomputes before every
consuming run and classifies drift as hard (shape/column changes, new nulls,
a column collapsing to constant: exit code 1, wired into any pre-run gate) or
soft (values changed, schema and stats hold). The discipline is two-phase and
cheap: fingerprint at build time, verify before spending GPU hours.

\subsection{Autonomous experimentation: write, smoke-test, dispatch, verify}
\label{sec:auto}

RIGOR's execution layer treats the agent as more than a dispatcher: the agent
\emph{writes and maintains the experiment code itself}, under the same
integrity gates as everything else. Given a training notebook, it applies the
four edits that make the notebook remotely runnable---a single tagged
parameter cell for injection~\cite{papermill}; an environment path shim so
the identical notebook runs on a laptop, on Kaggle, or on Colab; a loud
\texttt{SMOKE\_TEST} switch whose banner makes pipeline-check numbers
impossible to mistake for results; and a final cell exporting every metric,
type-cast, to a \texttt{results.json}---and then debugs its own work against
a local CPU environment. The \emph{smoke-test discipline} is a hard gate: no
configuration is dispatched to a GPU until the full pipeline has run
end-to-end on a tiny local subset, so debugging happens in seconds on the
laptop, never in minutes of donated GPU quota.

\textbf{run-remote} then drives the sweep unattended: parameter injection,
push to a free cloud GPU service, polling, artifact download, parsing, and an
append-only experiment journal. Failures are absorbed, not fatal: a failed
run is logged and retried under a fresh run identifier (so a stale journal
entry can never silently mask a retry) while the rest of the sweep proceeds.

The same layer scales to \emph{fleet dispatch}. A research group can register
each member's own cloud account and API token in a local roster (credentials
stay machine-level and never enter the repository), and the agent plans a
campaign across the pooled per-account quotas: it probes each account before
spending anything (a one-minute kernel confirming a GPU actually attaches),
stages datasets where each account needs them, partitions the sweeps, writes
one sweep file per account, launches the runners in parallel, monitors the
logs, self-heals transient failures, and returns a verified briefing rather
than a directory of raw files. The researcher's role collapses to two
moments: the kickoff instruction and the debrief. One hard rule survives
full autonomy: enabling GPU---which spends real, shared quota---requires the
owner's explicit approval, every time, in every project.

Version 1.2 adds \textbf{colab-run}, a second execution backend built around
an honest platform constraint: unlike Kaggle, free Google Colab exposes no
official headless-execution API, and RIGOR does not automate around a
provider's terms. The design is therefore \emph{semi-attended}: the agent
injects parameters (implementing the papermill injected-cell convention in
the standard library) and stages the notebook into the researcher's local
Google-Drive-synced folder, from which it appears in Colab within seconds;
execution is a single tap (\emph{Runtime}~$\to$~\emph{Run all}) in any
browser, including a phone; the notebook persists its \texttt{results.json}
back to Drive, and the agent's poller picks it up from the synced folder,
validates it, and journals it alongside the Kaggle runs. Everything but the
tap is automated---and the tap itself \emph{is} the explicit GPU approval
the safety rule requires, turning the human checkpoint from a procedural
promise into a structural necessity. The loop is \emph{field-verified}, not
merely designed: a staged notebook executed on a live Colab VM and returned
a token-stamped \texttt{results.json} through the synced folder, closing the
round trip end-to-end before the capability was claimed here
(Section~\ref{sec:case}).

\paragraph{Choosing a backend.} The two backends are not equals, and the
toolkit says so plainly. \textbf{Kaggle is the recommended default}: its
official kernels API makes the entire loop headless---push, poll,
download---which is what enables true fleet dispatch, overnight campaigns,
and the two-moment (kickoff/debrief) workflow. \textbf{colab-run is the
deliberate fallback}: one tap per run automates everything \emph{around}
execution but not execution itself, so a ten-run sweep means ten taps.
Reach for Colab when the week's Kaggle quota is exhausted, when a
Colab-specific runtime is wanted, or when the per-run human checkpoint is
itself desirable; route everything else through the Kaggle path.

\textbf{verify-run} is the checklist every completed run passes before its
numbers reach the researcher: all expected configurations present (a missing
key is a silent mid-run failure, not a partial success); smoke-test flags
honored (pipeline-check numbers are never reported as results); early-abort
(futility) stops disclosed rather than hidden; NaNs reported as failures;
parameter counts matched exactly against anchors (an exact match fingerprints
that the intended pipeline ran); and seed counts always disclosed, so a
single-seed number is never presented as a robust average.

\subsection{Continuity: the append-only lab notebook}
\label{sec:notebook}

Real investigations outlive any single agent session and fan out into parallel
tracks---a parsing track, a data-quality track, a modelling track---whose state
is exactly what conversation context loses: which track is blocked and why,
which result was superseded, what the next concrete action was.
\textbf{lab-notebook} makes a machine-readable notebook, not the conversation,
the source of truth. Entries are typed (\texttt{progress}, \texttt{finding},
\texttt{blocker}, \texttt{decision}, \texttt{correction}, \texttt{next}) and
appended to a per-investigation JSONL log that is never edited: a wrong entry
is superseded by a \texttt{correction} that points at it. A \texttt{finding} is
expected to cite the artifact that shows it (a CSV, a figure, a log)---the
script warns otherwise, and quoted numbers come from the artifact, not from
memory. Each session starts from a compact digest (every track's status and
last entry, open blockers with their reasons, queued next-steps with staleness
flags) and ends by recompiling \texttt{NOTEBOOK.md}, where any entry a
correction points at carries an explicit \emph{superseded-by} marker, so a
corrected number cannot be quoted unaware.

Two sub-agent workflows sit on top, and both keep the script as the truth
layer. \emph{Audit} spawns a read-only sub-agent that re-opens every finding's
evidence artifact and checks the quoted numbers against it---the script can
verify that evidence \emph{exists}; only an agent can verify the artifact still
\emph{shows the number}. Exercised live on a synthetic notebook, the auditing
sub-agent confirmed the seeded finding against its artifact and, unprompted,
flagged a numeric correction that carried no evidence of its own.
\emph{Narrate} addresses the single-coherent-story problem: a sub-agent given
\emph{only} the notebook writes the investigation as one narrative---the
hypothesis, what each track tried, dead ends and why they were abandoned, the
surviving evidence chain---in which every factual claim cites entry ids. The
fresh context is the point: it cannot remember the sessions, so if the story
cannot be written from the notebook alone, the notebook is incomplete, and the
sub-agent's gap list says exactly what to log. The \texttt{check-narrative}
subcommand then enforces the grounding mechanically: citations of nonexistent
entries fail, citations of superseded entries warn, and findings the narrative
never used are listed. One rule binds every sub-agent: entry identifiers are
assigned read-then-append, so sub-agents report back and only the main session
writes the log.

\subsection{The manuscript layer: local \LaTeX{}, agent-drawn figures, one-way numbers}
\label{sec:paper}

RIGOR's workflow keeps the entire manuscript---text, tables, and diagrams---as
plain text that compiles \emph{locally}, with the agent as typesetter and
illustrator. This deliberately displaces three tools researchers usually
reach for:

\begin{itemize}
\item \textbf{Overleaf} is replaced by a local \texttt{latexmk} build inside
the researcher's editor: no compile-minute limits, no internet dependence,
the full paper rebuilding in seconds, and the manuscript versioned in Git
alongside the code that produced its numbers.
\item \textbf{Diagramming tools (draw.io and kin)} are replaced by TikZ
\emph{written by the agent}: an architecture diagram is code, so the
researcher describes the figure in plain English and the agent writes the
drawing---an infinitely-zoomable vector figure that remains editable, by the
agent, for the life of the paper, instead of a mouse-drawn artifact that
must be redrawn by hand at every revision.
\item \textbf{AI image generators} are displaced for schematic figures by the
same mechanism: TikZ output is exact, editable, and diff-able, where
generated raster images are approximate, frozen, and blurry at print
resolution---and subject to usage caps besides.
\end{itemize}

What makes this workable in practice is that the agent \emph{closes its own
visual quality loop}: after compiling a figure it rasterizes the PDF, reads
the image back, spots overlapping labels, clipped boxes, or bad spacing,
edits the TikZ source, and recompiles---iterating until the figure is clean.
A loop that previously required a human, a mouse, and an afternoon runs
without either.

Result tables and result figures never pass through human hands at all: an
aggregator script reads every \texttt{results.json}, computes
mean~$\pm$~std across seeds, and writes \texttt{\textbackslash input}-ready
table fragments and figure PDFs directly into the manuscript tree. Numbers
therefore flow \emph{code $\rightarrow$ results.json $\rightarrow$
aggregator $\rightarrow$ generated tables and figures}, in one direction
only: no metric is ever retyped by hand into prose, and if a number looks
wrong the fix belongs in the code or the data, never in the text.
claims-audit (Section~\ref{sec:audit}) exists precisely to catch the places
where that rule was broken by hand before the pipeline enforced it.

\subsection{Closing the submission loop: the gate, the freeze, the rebuttal}
\label{sec:submission}

Two skills cover the lifecycle stages after the manuscript is written.
\textbf{submit-gate} runs the entire audit battery---bibliography, claims,
citations, statistics, data verification---as one command against a
per-project configuration, aggregating exit codes and output tails into a
single readiness report with a \textsc{ready}/\textsc{not-ready} verdict; a
failed required step is never argued down, only fixed or consciously waived
in the committed configuration. Immediately after submitting, \texttt{freeze}
snapshots (SHA-256) every file the submission's numbers rest on---tables,
results artifacts, the PDF---so that a reviewer question months later is
answered against \emph{what was submitted}, not against whatever the files
have since become; \texttt{verify-freeze} makes any divergence explicit, to
be disclosed in the response rather than silently papered over.

When reviews arrive, \textbf{rebuttal} turns them into a tracked checklist:
each comment gets a response with a declared action (\emph{change},
\emph{clarify}, or \emph{decline}), and every \emph{change} records the files
it claims to have modified and, optionally, a quote now present in the
manuscript. \texttt{check} then verifies each claim against the actual
revision diff---an anchor file absent from the diff, or a quote found
nowhere, fails the letter before a reviewer can catch the mismatch. The
classic rebuttal failure (``we have revised Section~3 accordingly,'' next to
a diff that shows nothing) becomes mechanically impossible to ship unnoticed.

\subsection{The reviewer's side: screening submitted PDFs}
\label{sec:review}

Everything above protects authors from their own tools. Two skills turn the
audits around so that the only input is a submitted paper's PDF---one paper,
or a conference batch delivered as an archive.

\textbf{ref-audit} extracts the reference list from the PDF text (the last
heading-like ``References'' onward), splits numbered entries by accepting a
bracket marker only when it continues the $1,2,3,\dots$ sequence---so inline
brackets inside an entry cannot derail the split---with an author-start
fallback for author--year styles, and parses each entry into title, venue,
year, DOI, arXiv identifier, and URL. Every entry is then verified through
bib-audit's engine, imported rather than forked, so authors and reviewers are
held to one verification policy. A printed arXiv identifier or DOI that
resolves nowhere is recorded as the strongest fabrication signal, but a DOI
damaged by PDF line-wrapping is first retried by title, so an extraction
artifact is never reported as a missing work. Per-entry checkpointing makes a
batch interruptible and resumable at the 1~request/s rate limit.

\textbf{result-audit} is the fully offline half and asks whether a paper
agrees with itself. It maps the extracted text to abstract, body, conclusion,
and back matter, and emits page-tagged, quoted candidates from seven checks:
headline numbers in the abstract or conclusion that sit close to a body or
table value but beyond its printed precision (\textsc{near-miss}) or appear
nowhere in the body (\textsc{orphan}); same-sentence ``improves by $N$\%''
claims that no pair of the sentence's numbers reproduces under either a
relative or a percentage-point reading; train/validation/test splits that do
not sum to 100; impossible statistics ($p>1$, accuracy above 100\%, $|r|>1$,
negative error metrics, and---with scipy present---statcheck-style
recomputation of $p$ from reported $t$, $F$, and $\chi^2$
statistics~\cite{nuijten2016statcheck}); means unattainable from the stated
integer $N$ (the GRIM test~\cite{brown2017grim}); and one metric quoted
irreconcilably in the abstract and the conclusion.

Both scripts stop where evidence stops. A \textsc{not-found} reference is a
\emph{candidate}: regional proceedings, non-English work, standards, theses,
and brand-new preprints are legitimate index misses. For each one, the agent
searches the \emph{cited} title and records \textsc{confirm} or
\textsc{exclude} with the evidence. A result candidate is checked against the
rendered PDF page and recorded as \textsc{confirmed-inconsistent} only with
both conflicting fragments quoted verbatim; otherwise it is \textsc{explained}
(both numbers are real but refer to different things) or
\textsc{extraction-noise}. A shared compiler merges both audits and the
adjudications into one triage workbook for the chair, one row per exact
finding---a named reference number, or the two conflicting quotes---in which
only adjudication-confirmed items become findings. Two rules are
non-negotiable. \emph{Confidentiality}: only the cited works' titles and
identifiers---metadata of already-published work---are sent to any API; the
submission's own title, authors, and text never leave the machine.
\emph{Facts, not attribution}: every output states what is verifiable (``the
printed arXiv identifier does not exist'') and nothing about how a paper was
written; judgments about AI use, and the decision, belong to the chair.

\section{Case study: a real project as first install}
\label{sec:case}

RIGOR's first install is a live research project---a multi-seed regional
electricity-load forecasting study with a LaTeX manuscript, a 14-entry
bibliography, a 245-paper literature collection, and six 10-seed experiment
studies executed remotely. Each tool produced concrete, verifiable findings:

\begin{itemize}
\item \textbf{bib-audit} (14 entries, 51\,s of wall time): one reference's title turned out to be
a paraphrase of the real paper's title (the entry resolved to nothing until
the audit surfaced the correctly-titled record); three title-only matches hit
different editions or records and were flagged with no auto-fix; two dataset
resources were verified by URL liveness; five missing DOIs were recovered.
The preprint-year calibration prevented three false alarms.
\item \textbf{claims-audit} (one manuscript): 271 numeric claims classified
(151 matched, 47 near-miss, 73 orphan) in 0.25\,s---mechanizing, and
extending, a manual audit that had found nine real text--data mismatches, and
leaving the human only the semantic adjudication of the flagged candidates.
\item \textbf{stat-check} (14 models $\times$ 10 seeds): the nominally best
model was statistically indistinguishable from several siblings (Wilcoxon
$p\approx0.5$), while every family-versus-baseline comparison survived exactly
($p=0.002$--$0.02$). A single-seed ``best architecture'' headline from an
earlier phase of the project did not survive; the published claim was reshaped
accordingly---the tool's purpose working as intended, against its own
project's preferred narrative.
\item \textbf{run-remote / verify-run}: the six 10-seed studies were executed
as unattended GPU sweeps---more than sixty runs plus retries, dispatched by
the agent across multiple registered free-tier accounts, every run passing
the verification checklist before aggregation. The pre-flight probe caught
one account whose GPU access was silently unavailable and rebalanced the
schedule onto the rest; a transient network outage mid-campaign was healed by
fresh-identifier retries, with no human intervention between kickoff and
briefing.
\item \textbf{colab-run} (field verification): the Colab loop was exercised
live before being claimed in this paper---a staged notebook executed on a
Colab VM and returned a token-stamped \texttt{results.json} through the
synced folder, matching the injected token and seed exactly. The field test
immediately exposed a real parameter-injection bug the unit tests had
missed: Python-style \texttt{True}/\texttt{False} values were being parsed
as JSON, injecting the truthy \emph{string} \texttt{'False'}, which would
have made smoke-test mode impossible to disable. It was fixed and locked in
as a regression test. We publish no capability we have not exercised---the
same standard the toolkit enforces on its users' claims.
\item \textbf{lit-review}: the collection was built from seven ranked queries,
snowballing, recommendations, and a boolean bulk pass; abstract enrichment
filled 9 of 39 missing abstracts from OpenAlex (the remainder are withheld by
publishers---an honest, documented gap); a single batch refresh re-dated all
citation counts. With the project's own description as the report's
\emph{focus}, the 245-paper collection tiers into 11 \textsc{core} / 95
\textsc{related} / 139 \textsc{peripheral}---and the \textsc{core} tier
surfaced, on its own, the same closest-rival paper a manual screen had
previously identified. Continued field use later surfaced two API edge-case
crashes---Semantic Scholar answering a literal \texttt{null} where a list was
expected, and a title whose comma collided with OpenAlex's reserved filter
syntax---both reproduced offline and locked in as regression tests, the same
path the parameter-injection bug above followed.
\item \textbf{manuscript layer}: all five of the case-study paper's
architecture and concept diagrams are agent-written TikZ, iterated under the
agent's own render-and-inspect loop; every result table and figure is
aggregator-generated from the raw run artifacts. Since the one-way pipeline
was adopted, no metric has been hand-typed into that manuscript's prose.
\item This paper's own bibliography was verified by resolving every DOI
through doi.org content negotiation before inclusion.
\end{itemize}

Nothing in the scripts is specific to the case study's domain. To make that
concrete---and reviewable without any API key, GPU, or agent
harness---the repository ships fully offline, synthetic worked examples
(\texttt{examples/}) for the statistics and audit tools, plus a network demo
in which \textbf{bib-audit} catches a deliberately fabricated citation; the
offline examples are exercised end-to-end by the CI test suite, so the
documented walkthroughs cannot drift from the code.

\section{Field test: a conference screening batch}
\label{sec:field}

The reviewer-side skills were run on a real screening batch before being
claimed here: 314 conference papers delivered as one archive. We report
aggregate counts only. No paper, author, venue, or individual finding is
identifiable from this section, and no submission content was sent to any
service.

\begin{table}[t]
\centering
\caption{Field test on a 314-paper conference screening batch (aggregate
counts only). Top: ref-audit verdicts. Bottom: result-audit candidates by
check and their adjudication outcome.}
\label{tab:field}
\small
\begin{tabular}{@{}lrr@{}}
\toprule
\multicolumn{3}{@{}l}{\textbf{ref-audit}: 5,294 entries from 303 completed papers}\\
\midrule
Verdict & Entries & Share\\
\midrule
\textsc{verified}      & 4,124 & 77.9\%\\
\textsc{mismatch}      &   725 & 13.7\%\\
\textsc{not-found}     &   293 &  5.5\%\\
\textsc{parse-failed}  &   136 &  2.6\%\\
\textsc{non-paper-ok}  &     9 &  0.2\%\\
\textsc{retracted}     &     4 &  0.1\%\\
\textsc{unverifiable}  &     3 &  0.1\%\\
\bottomrule
\end{tabular}

\vspace{0.8em}
\begin{tabular}{@{}lrrrr@{}}
\toprule
\multicolumn{5}{@{}l}{\textbf{result-audit}: 297 candidates, all adjudicated}\\
\midrule
Check & Candidates & Confirmed & Explained & Extraction noise\\
\midrule
\textsc{near-miss}   & 101 & 10 &  81 & 10\\
\textsc{orphan}      &  82 &  0 &  81 &  1\\
\textsc{split}       &  49 &  1 &  46 &  2\\
\textsc{arith}       &  47 &  0 &  46 &  1\\
\textsc{impossible}  &  11 &  0 &   0 & 11\\
\textsc{xsection}    &   7 &  4 &   3 &  0\\
\textsc{grim}        &   0 & -- &  -- & --\\
\midrule
Total                & 297 & 15 & 257 & 25\\
\bottomrule
\end{tabular}
\end{table}

\paragraph{References.} ref-audit completed 303 of the 314 papers (11 runs
were interrupted and had not been resumed at the time of writing), parsing
5,294 entries: 309 papers split in the numbered style, 5 via the author-start
fallback. Table~\ref{tab:field} (top) gives the verdicts. The 293
\textsc{not-found} entries were spread across 138 papers, and that spread is
the argument for mandatory adjudication. Of the 89 \textsc{not-found}
candidates adjudicated in the screening window, 58 were real works that the
bibliographic indexes do not cover, 6 stayed unresolved, and 25 had no trace
anywhere. Reported raw, the detector would have overstated the problem more
than threefold on the adjudicated set. Together with one \textsc{mismatch}
the adjudicator confirmed as untraceable, 26 references across 5 papers
reached the chair as confirmed findings. The four \textsc{retracted} verdicts
come mechanically from Retraction Watch data. \textsc{mismatch} entries
(13.7\%) were not adjudicated in this pass; they mix wrong-edition matches
with misquoted titles and years, and are reported as drift, not fabrication.
At the rate limit this is the slow half of the screening: its checkpoints span
about 27 hours of wall-clock time across resumed sessions, pauses included.

\paragraph{Internal consistency.} result-audit processed all 314 papers in
under two minutes. It checked 1,312 headline claims from abstracts and
conclusions, of which 1,129 (86\%) matched a body or table value at printed
precision, and raised 297 candidates in 116 papers. Every candidate was
adjudicated against the rendered page (Table~\ref{tab:field}, bottom): 15
were confirmed inconsistencies across 9 papers---10 near-misses, 4
abstract-versus-conclusion contradictions, and 1 split that does not sum to
100. The per-check yield is worth stating plainly, because it tells a user
which flags to read first. The abstract-versus-conclusion check was the most
precise (4 of 7 confirmed) and near-miss confirmed about one candidate in ten.
The arithmetic, orphan, and split checks confirmed almost nothing: their
candidates were overwhelmingly real numbers that referred to a different
baseline, dataset, or metric variant. All 11 impossible-statistic flags were
artifacts of PDF text extraction. Overall, 5\% of candidates survived
adjudication. That is the designed trade of false positives for coverage, and
the reason the chair's workbook admits only adjudicated findings.

\section{Limitations}

RIGOR's grounding is only as complete as its substrate: Semantic Scholar and
OpenAlex have coverage gaps (some publishers withhold abstracts; very recent
or non-English work may be missing), so ``not found'' is evidence, not proof,
of non-existence. Rate limits (1~request/s on the primary API) make large
collections slow to build and briefly throttle bursty use.
The reviewer-side skills inherit that coverage gap with higher stakes: a
\textsc{not-found} is likelier for regional, non-English, and grey-literature
venues, so raw counts would disproportionately flag authors from
under-indexed communities---which is why only adjudicated findings reach a
chair, and why no output says anything about how a paper was written. PDF
text extraction is lossy (two-column reflow, mangled tables), and the field
test's per-check yields (Section~\ref{sec:field}) show that some checks are
precise while others are coverage nets whose flags are mostly explained.
The claims auditor is calibrated for LaTeX manuscripts with generated tables;
purely hand-written papers reduce it to placeholder and figure checks. The
statistical tool assumes per-seed results files in a simple JSON schema;
projects must export it (a few lines in the training notebook). Paired tests
at $n{\approx}10$ have limited power, which the tool discloses rather than
solves. Finally, the skill layer's semantic adjudication depends on the
hosting agent following written rules; the guarantees that matter most
(no citation outside the store, no auto-edits, exact $p$-values) are enforced
by the tools themselves, but the quality of, e.g., near-miss adjudication
varies with the agent.

\section{Conclusion}

RIGOR treats research integrity as an architecture problem: identify the
places where AI assistance can silently fabricate, drift, or overclaim, and
put a mechanical, auditable gate at each one. The result is a small toolkit
that a researcher can copy into any project, drive with or without an agent,
and extend---with the guarantee that everything it lets through leaves a
verifiable trail. The same gates serve the reviewer as well as the author.

\paragraph{Future work.} The execution layer is backend-agnostic by design
(prepare, stage, execute, collect, verify), and v1.2's Colab backend already
lets groups pool quota across two free services. Next: fully headless Colab
execution through the paid Colab Enterprise API for funded groups, scheduled
topic watching for teams with always-on machines, a multi-user
deployment study of RIGOR across a working research group, and, for the
reviewer-side skills, a multi-committee deployment that measures agreement
between independent adjudicators.

\paragraph{Availability.} RIGOR is available under the MIT license at
\url{https://github.com/alif199339/rigor}. Collections built with it
rely on the Semantic Scholar API; published material based on them should cite
Kinney et al.~\cite{kinney2023semanticscholar} per the API license.

\bibliographystyle{unsrt}
\bibliography{refs}

\end{document}
